% =====================================================================
%  Lecture 07 -- LTE Scheduling                           (Persian edition)
% =====================================================================
\documentclass[11pt,a4paper]{article}
\providecommand\ACNroot{../..}
\input{\ACNroot/templates/lecture-fa}

\ACNset{number}{07}
\ACNset{title}{زمان‌بندی در \lr{LTE}}
\ACNset{shorttitle}{زمان‌بندی در \lr{LTE}}
\ACNset{titleother}{LTE Scheduling}
\ACNset{part}{بخش دوم -- دسترسی و هستهٔ نسل چهارم (درس‌های ۵ تا ۱۳)}
\ACNset{time}{۸ تا ۱۰ ساعت (۳٫۵ ساعت مطالعه، ۴ تا ۶ ساعت تمرین همراه با شبیه‌سازی)}
\ACNset{level}{متوسط تا پیشرفته}
\ACNset{prereq}{درس‌های ۵ و ۶؛ ظرفیت شانون؛ محوشدگی رایلی؛ بهینه‌سازی پایه}
\ACNset{gen}{4}

\begin{document}
\ACNcover

% ---------------------------------------------------------------------
\section{اهداف یادگیری}
پس از مطالعهٔ این درس باید بتوانید:
\begin{objectives}
  \item شبکهٔ منابع \acro{LTE} -- قاب، زیرقاب، شیار، عنصر منبع، بلوک منبع -- را \textbf{توصیف کنید} و منابع
        هر پهنای باند \acro{LTE} را \textbf{بشمارید}.
  \item سیگنالینگ زمان‌بندی فروسو و فراسو را \textbf{توضیح دهید}: \lr{DCI}، مجوزها، \lr{BSR}، \lr{PHR}، \lr{CQI}.
  \item تطبیق پیوند را \textbf{به کار ببرید}: \acro{CQI} را به مدولاسیون و کدگذاری نگاشت کنید و تطبیق پیوند
        حلقهٔ بیرونی (\acro{OLLA}) را \textbf{تحلیل کنید}.
  \item \acro{HARQ} با هشت فرایند توقف-و-انتظار را \textbf{تحلیل کنید} و اثر آن بر گذردهی را محاسبه کنید.
  \item زمان‌بندی نوبتی، بیشینهٔ \lr{C/I} و منصفانهٔ تناسبی را \textbf{مقایسه کنید} و \textbf{استخراج کنید} چرا
        \acro{PF} مجموع لگاریتم گذردهی‌ها را بیشینه می‌کند.
  \item مصالحهٔ گذردهی--انصاف--کیفیت خدمات را با شاخص \lr{Jain} و یک شبیه‌سازی \textbf{ارزیابی کنید}.
\end{objectives}

% ---------------------------------------------------------------------
\section{دانش پیش‌نیاز}
به پشتهٔ پروتکل \acro{LTE} (درس ۶)، به‌ویژه نقش‌های زمان‌بندی و \acro{HARQ} لایهٔ \acro{MAC}، ظرفیت شانون
$C = B\log_2(1+\SNR)$ و این واقعیت نیاز دارید که در محوشدگی رایلی \acro{SNR} لحظه‌ای حول میانگین خود توزیع
نمایی دارد. حسابان پایه (مشتق، استدلال‌های شبیه لاگرانژ) در \cref{sec:pf} به کار می‌رود.

\paragraph{واژه‌های کلیدی} واژگان این درس همراه با معادل انگلیسی آن‌ها:
\begin{keyterms}
عنصر منبع & \lr{Resource element (RE)} & یک زیرحامل در طول یک نماد \acro{OFDM}. \\
بلوک منبع & \lr{Resource block (RB)} & ۱۲ زیرحامل (\lr{180\,kHz}) $\times$ یک شیار. \\
بازهٔ زمانی ارسال & \lr{TTI} & بازهٔ زمان‌بندی: یک زیرقاب ۱ میلی‌ثانیه‌ای در \acro{LTE}. \\
اطلاعات کنترلی فروسو & \lr{DCI} & پیام زمان‌بندی روی \acro{PDCCH}: منابع، \acro{MCS}، \acro{HARQ}. \\
نشانگر کیفیت کانال & \lr{CQI} & گزارش \acro{UE}: بالاترین \acro{MCS} قابل آشکارسازی با \lr{BLER} ده درصد. \\
طرح مدولاسیون و کدگذاری & \lr{MCS} & مرتبهٔ مدولاسیون + نرخ کد که زمان‌بند انتخاب می‌کند. \\
گزارش وضعیت بافر / حاشیهٔ توان & \lr{BSR / PHR} & گزارش‌های فراسو: دادهٔ در صف / توان باقی‌مانده. \\
درخواست تکرار خودکار ترکیبی & \lr{HARQ} & ارسال مجدد سریع با ترکیب نرم. \\
منصفانهٔ تناسبی & \lr{Proportional fair (PF)} & زمان‌بندی که $\sum\log R_k$ را بیشینه می‌کند. \\
تنوع چندکاربره & \lr{Multi-user diversity} & بهره از خدمت‌دهی به کاربران در قله‌های محوشدگی‌شان. \\
\end{keyterms}

% ---------------------------------------------------------------------
\section{مقدمه و انگیزه}
هر میلی‌ثانیه، \acro{eNB} در \acro{LTE} باید به یک پرسش پاسخ دهد: \emph{چه کسی کدام بلوک‌های منبع را با کدام
مدولاسیون و کدگذاری می‌گیرد؟} این پاسخ همهٔ آنچه کاربران درک می‌کنند -- گذردهی، تأخیر، انصاف -- و همهٔ آنچه
اپراتور هزینه‌اش را می‌پردازد تعیین می‌کند، زیرا طیف گران‌ترین منبع یک شبکهٔ موبایل است.

مسئله دشوار است، زیرا کانال هر کاربر نوسان می‌کند (محوشدگی سریع، تحرک، تداخل)، ترافیک انفجاری است و اهداف
با هم تعارض دارند: بیشینه کردن گذردهی کل به نفع کاربران نزدیک ایستگاه پایه است، در حالی که انصاف و تضمین
کیفیت خدمات به خدمت‌دهی به کاربران لبهٔ سلول نیاز دارد. \acro{LTE} ابزاری ریزدانه -- ۱ میلی‌ثانیه در زمان و
\lr{180\,kHz} در فرکانس -- به \acro{eNB} می‌دهد و \emph{الگوریتم} زمان‌بندی را به سازنده می‌سپارد. \acro{3GPP}
سیگنالینگ (آنچه \acro{UE} گزارش می‌دهد و نحوهٔ ارسال مجوزها) را استاندارد می‌کند، نه تصمیم را. این درس هر دو
را پوشش می‌دهد: سازوکارهای استاندارد و الگوریتم‌های کلاسیک همراه با مبانی ریاضی آن‌ها.

% ---------------------------------------------------------------------
\section{مفاهیم بنیادی}

\subsection{شبکهٔ منابع}
در حوزهٔ زمان، یک قاب رادیویی \acro{LTE} ده میلی‌ثانیه طول دارد و شامل ده \emph{زیرقاب} یک‌میلی‌ثانیه‌ای است
که هر یک از دو \emph{شیار} ۰٫۵ میلی‌ثانیه‌ای با ۷ نماد \acro{OFDM} (پیشوند چرخه‌ای عادی) تشکیل شده است
\cite{ts36211}. در حوزهٔ فرکانس فاصلهٔ زیرحامل‌ها $\Delta f = 15$\,\lr{kHz} است. یک \emph{عنصر منبع} یک
زیرحامل برای یک نماد است؛ یک \emph{بلوک منبع} ۱۲ زیرحامل (\lr{180\,kHz}) برای یک شیار (۸۴ عنصر منبع) است.
زمان‌بند معمولاً در هر \acro{TTI} یک‌میلی‌ثانیه‌ای \emph{جفت \lr{RB}} (هر دو شیار یک زیرقاب، ۱۶۸ عنصر منبع)
تخصیص می‌دهد.

\begin{table}[htbp]
\centering\small
\caption{پهنای باند کانال‌های \acro{LTE} و بلوک‌های منبع \cite{ts36211}.}\label{tab:bw}
\begin{tabular}{@{}lcccccc@{}}
\toprule
پهنای باند کانال [\lr{MHz}] & 1.4 & 3 & 5 & 10 & 15 & 20 \\
بلوک‌های منبع $N_{\mathrm{RB}}$ & 6 & 15 & 25 & 50 & 75 & 100 \\
پهنای باند اشغال‌شده [\lr{MHz}] & 1.08 & 2.7 & 4.5 & 9.0 & 13.5 & 18.0 \\
\bottomrule
\end{tabular}
\end{table}

\begin{acnfigure}
  \resizebox{\linewidth}{!}{\acnfig{lecture-07/resource-grid}}
  \caption{یک جفت \acro{RB} فروسو (۱۲ زیرحامل $\times$ ۱۴ نماد) با ناحیهٔ کنترل دونمادی و سیگنال‌های مرجع
  ویژهٔ سلول دو درگاهی.}
  \label{fig:grid}
\end{acnfigure}

\subsection{زمان‌بند چه می‌داند}
\begin{itemize}
  \item \textbf{کیفیت کانال}: گزارش‌های \acro{CQI} (پهن‌باند یا برای هر زیرباند)، نشانگرهای رتبه و پیش‌کدگذار
        (\lr{RI, PMI}) در فروسو؛ سیگنال‌های مرجع کاوشی (\acro{SRS}) در فراسو.
  \item \textbf{ترافیک}: بافرهای فروسو در \acro{eNB} هستند؛ بافرهای فراسو را \acro{UE} در \emph{گزارش‌های وضعیت
        بافر} (\acro{BSR})، برای هر گروه از کانال‌های منطقی، گزارش می‌دهد.
  \item \textbf{کیفیت خدمات}: \acro{QCI}، نرخ بیت تضمین‌شده (\acro{GBR}) و اولویت هر حامل (درس ۲۲).
  \item \textbf{توان \acro{UE}}: \emph{گزارش‌های حاشیهٔ توان} (\acro{PHR}) می‌گویند \acro{UE} چه مقدار توان بیشتر
        می‌تواند استفاده کند تا زمان‌بند فراسو بیش از توان \acro{UE} \acro{RB} تخصیص ندهد.
\end{itemize}

% ---------------------------------------------------------------------
\section{توضیحات فنی تفصیلی}

\subsection{زمان‌بندی فروسو}
برای هر \acro{UE} زمان‌بندی‌شده، \acro{eNB} یک پیام \emph{اطلاعات کنترلی فروسو} (\acro{DCI}) را روی \acro{PDCCH}
در ناحیهٔ کنترل همان زیرقاب می‌فرستد \cite{ts36213}. یک \acro{DCI} شامل تخصیص منابع، شاخص \acro{MCS}، شمارهٔ
فرایند \acro{HARQ}، نشانگر دادهٔ جدید، نسخهٔ افزونگی و یک فرمان کنترل توان برای \acro{PUCCH} است. سه نوع تخصیص
منابع اندازهٔ سیگنالینگ را با انعطاف معاوضه می‌کنند: نوع ۰ (بیت‌نقشهٔ گروه‌های \acro{RB})، نوع ۱ (\acro{RB}ها در
یک زیرمجموعه) و نوع ۲ (مجموعه‌ای پیوسته، فشرده). \acro{UE} کانال \acro{PDSCH} را در همان زیرقاب آشکار می‌کند و
چهار زیرقاب بعد \lr{ACK/NACK} برمی‌گرداند.

\subsection{زمان‌بندی فراسو}
فراسو از سه جهت متفاوت است. (۱) \acro{eNB} بافر \acro{UE} را نمی‌بیند: \acro{UE} یک \emph{درخواست زمان‌بندی}
(\acro{SR}) روی \acro{PUCCH} می‌فرستد، مجوز کوچکی می‌گیرد و سپس بافر خود را با یک \acro{BSR} گزارش می‌دهد.
(۲) مجوزها چهار زیرقاب پیش از ارسال \acro{PUSCH} روی \acro{PDCCH} فرستاده می‌شوند (\acro{FDD}). (۳) \acro{SC-FDMA}
ایجاب می‌کند \acro{RB}های یک \acro{UE} در فرکانس \emph{پیوسته} باشند تا نسبت توان قله به میانگین پایین بماند، که
زمان‌بندی گزینشی در فرکانس را محدود می‌کند. \acro{HARQ} فراسو \emph{همگام} است: ارسال‌های مجدد دقیقاً ۸ میلی‌ثانیه
بعد رخ می‌دهند و روی \acro{PHICH} تأیید می‌شوند.

\subsection{\lr{CQI} و تطبیق پیوند}
\acro{UE} یک \acro{CQI} چهاربیتی گزارش می‌دهد: بالاترین شاخصی که مدولاسیون و نرخ کد آن با احتمال خطای بلوک
انتقال حداکثر ۱۰٪ دریافت شود \cite{ts36213}. \Cref{tab:cqi} جدول استاندارد را نشان می‌دهد. \acro{eNB} مقدار
\acro{CQI} را (پس از اصلاح، \acro{OLLA} را در ادامه ببینید) به یک شاخص \acro{MCS} و اندازهٔ بلوک انتقال نگاشت می‌کند.

\begin{table}[htbp]
\centering\small
\caption{جدول \acro{CQI} چهاربیتی \acro{LTE} (\lr{TS 36.213, Table 7.2.3-1}). بازده بر حسب بیت اطلاعات به ازای
هر عنصر منبع است.}\label{tab:cqi}
\rowcolors{2}{acnmist}{white}
\begin{tabular}{@{}cccc|cccc@{}}
\toprule
\rowcolor{white}\acro{CQI} & مدولاسیون & نرخ کد$\times1024$ & بازده & \acro{CQI} & مدولاسیون & نرخ کد$\times1024$ & بازده \\
\midrule
1 & \lr{QPSK} & 78  & 0.1523 & 9  & \lr{16QAM} & 616 & 2.4063 \\
2 & \lr{QPSK} & 120 & 0.2344 & 10 & \lr{64QAM} & 466 & 2.7305 \\
3 & \lr{QPSK} & 193 & 0.3770 & 11 & \lr{64QAM} & 567 & 3.3223 \\
4 & \lr{QPSK} & 308 & 0.6016 & 12 & \lr{64QAM} & 666 & 3.9023 \\
5 & \lr{QPSK} & 449 & 0.8770 & 13 & \lr{64QAM} & 772 & 4.5234 \\
6 & \lr{QPSK} & 602 & 1.1758 & 14 & \lr{64QAM} & 873 & 5.1152 \\
7 & \lr{16QAM} & 378 & 1.4766 & 15 & \lr{64QAM} & 948 & 5.5547 \\
8 & \lr{16QAM} & 490 & 1.9141 & 0  & \multicolumn{3}{l}{خارج از محدوده} \\
\bottomrule
\end{tabular}
\end{table}

توجه کنید که بازده صرفاً $Q_m \times$ (نرخ کد) است: برای \acro{CQI} ۱۵، $6\times948/1024 = 5.5547$ بیت به ازای هر
عنصر منبع. \acro{CQI} پهن‌باند کل حامل را توصیف می‌کند؛ \acro{CQI} زیرباند (برای هر گروه از \acro{RB}ها) زمان‌بندی
گزینشی در فرکانس را ممکن می‌کند اما هزینهٔ سربار فراسو دارد.

\subsection{\lr{HARQ}}
\acro{LTE} در \acro{FDD} از هشت فرایند موازی توقف-و-انتظار \acro{HARQ} استفاده می‌کند (\cref{fig:harq}): در حالی که
فرایند ۲ منتظر تأیید خود است، فرایندهای ۳ تا ۸ از زیرقاب‌های بعدی استفاده می‌کنند و کانال هرگز بیکار نمی‌ماند. با
\emph{ترکیب چیس}، ارسال مجدد همان بیت‌های کدشده را تکرار می‌کند و گیرنده انرژی‌های سیگنال را جمع می‌کند؛ با
\emph{افزونگی افزایشی}، بیت‌های توازن جدیدی ارسال می‌شود که نرخ کد مؤثر را کاهش می‌دهد. عملکرد معمول
\lr{BLER} ارسال نخست ۱۰٪ و \lr{BLER} باقی‌مانده پس از \acro{HARQ} در حدود $10^{-3}$ یا کمتر را هدف می‌گیرد؛
\acro{ARQ} در \acro{RLC} (درس ۶) بقیه را می‌گیرد.

\begin{acnfigure}
  \acnfig{lecture-07/harq-processes}
  \caption{هشت فرایند \acro{HARQ} در \acro{FDD}. یک \acro{NACK} برای فرایند \lr{P2} چهار زیرقاب پس از ارسال
  می‌رسد؛ ارسال مجدد هشت زیرقاب پس از ارسال اصلی رخ می‌دهد.}
  \label{fig:harq}
\end{acnfigure}

\subsection{الگوریتم‌های زمان‌بندی}
\paragraph{نوبتی (\lr{RR})} کاربران بدون توجه به کیفیت کانال به نوبت خدمت می‌گیرند. در \emph{زمان} کاملاً
منصفانه است، اما کانال را نادیده می‌گیرد و قله‌های خوب محوشدگی را هدر می‌دهد.
\paragraph{بیشینهٔ \lr{C/I} (بیشینهٔ نرخ)} در هر \acro{TTI} و هر \acro{RB} به کاربری با بالاترین نرخ لحظه‌ای
$r_k(t)$ خدمت می‌دهد. گذردهی سلول را بیشینه می‌کند اما کاربران لبهٔ سلول را گرسنه می‌گذارد.
\paragraph{منصفانهٔ تناسبی (\lr{PF})} به کاربری خدمت می‌دهد که بزرگ‌ترین نسبت نرخ لحظه‌ای به میانگین گذردهی
خودش را دارد \cite{jalali2000hdr}:
\begin{equation}\label{eq:pfmetric}
  k^\ast(t) = \argmax_k \frac{r_k(t)}{\bar R_k(t)}, \qquad
  \bar R_k(t+1) = \Bigl(1-\tfrac{1}{t_\mathrm{c}}\Bigr)\bar R_k(t) + \tfrac{1}{t_\mathrm{c}}\,r_k(t)\,\mathbb{1}\{k = k^\ast(t)\},
\end{equation}
که $t_\mathrm{c}$ پنجرهٔ میانگین‌گیری بر حسب \acro{TTI} است. هر کاربر وقتی خدمت می‌گیرد که کانالش \emph{نسبت به
میانگین خودش} خوب باشد؛ پس کاربران لبهٔ سلول نیز قله‌های خود را می‌گیرند.
\paragraph{زمان‌بندی آگاه از کیفیت خدمات} زمان‌بندهای واقعی ابتدا حامل‌های \acro{GBR} را برآورده می‌کنند (مثلاً
\acro{VoLTE}، اغلب با زمان‌بندی نیمه‌پایدار)، معیارها را با اولویت \acro{QCI} و بودجهٔ تأخیر بسته وزن می‌دهند و
برای ترافیک بهترین-تلاش باقی‌مانده از قواعد شبیه \acro{PF} استفاده می‌کنند.

% ---------------------------------------------------------------------
\section{مدل‌های ریاضی}

\subsection{از \lr{CQI} تا نرخ داده}\label{sec:rate}
برای \acro{UE}ای که $N_{\mathrm{RB}}$ جفت \acro{RB} با $N_{\mathrm{RE}}^{\mathrm{data}}$ عنصر منبع داده در هر کدام و
بازده $\eta_{\mathrm{CQI}}$ (بیت بر عنصر منبع) از \cref{tab:cqi} گرفته است، نرخ برای هر لایه برابر است با
\begin{equation}\label{eq:cqirate}
  R = \eta_{\mathrm{CQI}}\,\frac{N_{\mathrm{RB}}\,N_{\mathrm{RE}}^{\mathrm{data}}}{T_{\mathrm{sf}}}.
\end{equation}
\emph{مثال:} \acro{CQI} ۹ ($\eta = 2.4063$)، ۵۰ \acro{RB}، ۱۳۲ عنصر منبع داده (دو نماد کنترل، دو درگاه \acro{CRS}):
$R = 2.4063\times50\times132/10^{-3} = \qty{15.9}{\mega\bit\per\second}$.

\subsection{تطبیق پیوند حلقهٔ بیرونی}\label{sec:olla}
\paragraph{مسئله} گزارش‌های \acro{CQI} تأخیردار و اریب‌اند (خطاهای تخمین، تغییر تداخل)؛ اگر \acro{eNB} مستقیماً از
آن‌ها استفاده می‌کرد، \lr{BLER} واقعی از هدف ۱۰٪ فاصله می‌گرفت.
\paragraph{مدل} \acro{eNB} پیش از انتخاب \acro{MCS} آفست $\Delta$ (دسی‌بل) را به \acro{SINR} گزارش‌شده می‌افزاید. پس
از هر \acro{ACK} مقدار $\Delta$ را $\delta_{\mathrm{up}}$ بالا و پس از هر \acro{NACK} آن را $\delta_{\mathrm{down}}$
پایین می‌برد.
\paragraph{استخراج} در حالت پایدار تغییر مورد انتظار صفر است. اگر $p$ مقدار \lr{BLER} در نقطهٔ کار باشد،
\begin{equation}\label{eq:olla}
  (1-p)\,\delta_{\mathrm{up}} - p\,\delta_{\mathrm{down}} = 0
  \quad\Longrightarrow\quad
  p = \frac{\delta_{\mathrm{up}}}{\delta_{\mathrm{up}}+\delta_{\mathrm{down}}}.
\end{equation}
بنابراین انتخاب $\delta_{\mathrm{up}} = \delta_{\mathrm{down}}\,p_\mathrm{t}/(1-p_\mathrm{t})$ صرف‌نظر از اریبی \acro{CQI}،
\lr{BLER} را به هدف $p_\mathrm{t}$ می‌راند. برای $p_\mathrm{t} = 0.1$ و $\delta_{\mathrm{down}} = 0.5$\,\lr{dB}، مقدار
$\delta_{\mathrm{up}} = 0.056$\,\lr{dB} است.
\paragraph{تفسیر} \acro{OLLA} یک حلقهٔ کنترل تصادفی است: گام‌های کوچک \lr{BLER} دقیق اما همگرایی کند (حدود
$1/\delta$ گام) می‌دهند و گام‌های بزرگ سریع واکنش نشان می‌دهند اما \acro{MCS} را نوسانی می‌کنند.

\subsection{گذردهی \lr{HARQ}}
اگر ارسال نخست با احتمال $p_1$ و ارسال مجدد (پس از ترکیب) با احتمال $p_2 < p_1$ شکست بخورد، میانگین تعداد
ارسال‌ها به ازای هر بلوک انتقال (با چشم‌پوشی از شکست‌های تلاش سوم) $\bar N \approx 1 + p_1 + p_1p_2$ و بازده
گذردهی $1/\bar N$ است. با $p_1 = 0.1$ و $p_2 = 0.01$: $\bar N = 1.101$ و بازده ۹۰٫۸٪. هدف \lr{BLER} ده درصد تنها
حدود ۹٪ گذردهی هزینه دارد اما \acro{MCS} بسیار بالاتری از هدف ۱٪ را ممکن می‌کند -- دلیل اینکه \acro{LTE} در ۱۰٪
کار می‌کند.

\subsection{چرا منصفانهٔ تناسبی «به‌طور تناسبی منصفانه» است}\label{sec:pf}
\paragraph{ادعا} در میان همهٔ سیاست‌های زمان‌بندی، \acro{PF} (به‌طور مجانبی و برای پنجره‌های میانگین‌گیری طولانی)
تابع سودمندی زیر را بیشینه می‌کند:
\begin{equation}\label{eq:util}
  U = \sum_{k=1}^{K}\log \bar R_k .
\end{equation}
\paragraph{استدلال} دادن \acro{TTI} فعلی به کاربر $k$ را در نظر بگیرید. در تقریب مرتبهٔ اول، $\bar R_k$ را به اندازهٔ
$\epsilon\,r_k(t)$ (با $\epsilon = 1/t_\mathrm{c}$) افزایش می‌دهد و سودمندی را به این اندازه تغییر می‌دهد:
\begin{equation}
  \Delta U \approx \frac{\partial U}{\partial \bar R_k}\,\epsilon\,r_k(t) = \epsilon\,\frac{r_k(t)}{\bar R_k}.
\end{equation}
بنابراین انتخاب حریصانه‌ای که در هر \acro{TTI} افزایش سودمندی را بیشینه می‌کند دقیقاً $\argmax_k r_k(t)/\bar R_k$،
یعنی \cref{eq:pfmetric}، است. چون $U$ مقعر است، گام‌های حریصانه به بیشینهٔ آن همگرا می‌شوند. یک ویژگی سودمند
نقطهٔ بهینه \emph{انصاف تناسبی} است: هر تغییر شدنی در بردار گذردهی رابطهٔ $\sum_k \Delta\bar R_k/\bar R_k \le 0$
را برقرار می‌کند -- هیچ کاربری نمی‌تواند درصدی از گذردهی را به دست آورد بی‌آنکه دیگران دست‌کم همان درصد کل را از
دست بدهند. بیشینهٔ \lr{C/I} به‌جای آن $\sum \bar R_k$ را بیشینه می‌کند و \acro{RR} سهم‌های زمانی را برابر می‌کند.

\subsection{شاخص انصاف \lr{Jain}}
برای مقایسهٔ زمان‌بندها از این شاخص استفاده می‌کنیم \cite{jain1984fairness}:
\begin{equation}\label{eq:jain}
  J = \frac{\bigl(\sum_{k=1}^{K}\bar R_k\bigr)^2}{K\sum_{k=1}^{K}\bar R_k^2}, \qquad \frac{1}{K}\le J\le 1,
\end{equation}
که وقتی همهٔ کاربران گذردهی یکسان دارند برابر ۱ و وقتی یک کاربر همه‌چیز را می‌گیرد برابر $1/K$ است.

\subsection{شبیه‌سازی: تنوع چندکاربره}\label{sec:sim}
\paragraph{تنظیمات} \lr{\texttt{scripts/scheduler\_sim.py}} یک حامل ده‌مگاهرتزی (۹ مگاهرتز اشغال‌شده) را شبیه‌سازی
می‌کند که در هر \acro{TTI} به‌طور کامل به یک کاربر داده می‌شود؛ چهار کاربر با میانگین \acro{SNR} برابر ۲۰، ۱۲، ۵
و ۰ دسی‌بل؛ محوشدگی بلوکی رایلی مستقل در هر \acro{TTI}؛ بازده $\min(0.75\log_2(1+\SNR),5.55)$ بیت بر ثانیه بر
هرتز (نگاشت شانون تضعیف‌شده با سقف \acro{CQI} ۱۵)؛ ۲۰۰۰۰ \acro{TTI}؛ $t_\mathrm{c} = 100$؛ بذر تصادفی ثابت.
\paragraph{نتایج} (\cref{fig:sched}، \cref{tab:sim})
\begin{table}[htbp]
\centering\small
\caption{گذردهی شبیه‌سازی‌شدهٔ هر کاربر و سلول [\lr{Mbit/s}] و شاخص \lr{Jain}.}\label{tab:sim}
\begin{tabular}{@{}lrrrrrr@{}}
\toprule
زمان‌بند & \lr{UE1 (20\,dB)} & \lr{UE2 (12\,dB)} & \lr{UE3 (5\,dB)} & \lr{UE4 (0\,dB)} & سلول & $J$ \\
\midrule
نوبتی           & 9.74  & 5.83 & 2.88 & 1.43 & 19.88 & 0.710 \\
منصفانهٔ تناسبی & 12.97 & 8.34 & 4.60 & 2.50 & 28.40 & 0.761 \\
بیشینهٔ \lr{C/I} & 35.73 & 4.18 & 0.21 & 0.01 & 40.13 & 0.311 \\
\bottomrule
\end{tabular}
\end{table}
\paragraph{تفسیر} \acro{PF} به \emph{هر} کاربر گذردهی بیشتری از \acro{RR} می‌دهد: هر کاربر هنوز حدود یک‌چهارم
\acro{TTI}ها را می‌گیرد، اما این‌ها بهترین \acro{TTI}های اوست. این همان \emph{تنوع چندکاربره} است، بهره‌ای که با
تعداد کاربران و عمق محوشدگی رشد می‌کند. بیشینهٔ \lr{C/I} مجموع سلول را ۴۱٪ دیگر بالا می‌برد اما عملاً \lr{UE3} و
\lr{UE4} را خاموش می‌کند. مدل گزینش‌پذیری فرکانسی (زمان‌بندی زیرباند تنوع بیشتری می‌افزاید)، تأخیر و خطای
\acro{CQI} و پویایی ترافیک را نادیده می‌گیرد (همهٔ کاربران همیشه داده دارند -- فرض «بافر پر»).

\begin{acnfigure}
  \acnfig{lecture-07/scheduler-results}
  \caption{گذردهی هر کاربر با زمان‌بندی نوبتی، منصفانهٔ تناسبی و بیشینهٔ \lr{C/I} (شبیه‌سازی \cref{sec:sim}).}
  \label{fig:sched}
\end{acnfigure}

% ---------------------------------------------------------------------
\section{معماری شبکه و نمودارها}
زمان‌بندی کاملاً کارکرد \acro{MAC} در \acro{eNB} است؛ تنها ورودی‌های سطح شبکه پارامترهای کیفیت خدمات هر حامل
هستند که هسته هنگام برقراری حامل‌ها از طریق \lr{S1AP} تعیین می‌کند (\lr{QCI, ARP, GBR, MBR, AMBR}) (درس‌های ۹ و
۲۲). میان \acro{eNB}ها، پیام‌های \emph{اطلاعات بار} \acro{X2} هماهنگی تداخل میان سلولی را ممکن می‌کنند: همسایه‌ها
می‌توانند توافق کنند که کاربران لبهٔ سلول را روی \acro{RB}های یکسان زمان‌بندی نکنند. \Cref{fig:trilemma} الگوریتم‌ها
را در فضای اهداف جای می‌دهد.

\begin{acnfigure}
  \acnfig{lecture-07/trilemma}
  \caption{سه‌راهی زمان‌بندی. هیچ الگوریتمی گذردهی، انصاف و کیفیت خدمات را هم‌زمان بیشینه نمی‌کند؛ زمان‌بندهای
  عملی یک مرحلهٔ کیفیت خدمات را با یک مرحلهٔ شبیه \acro{PF} ترکیب می‌کنند.}
  \label{fig:trilemma}
\end{acnfigure}

% ---------------------------------------------------------------------
\section{رویه‌های گام‌به‌گام}
\begin{procedurebox}{یک \lr{TTI} از زمان‌بندی فروسو}\label{proc:dl}
\begin{steps}
  \item گردآوری ورودی‌ها: آخرین \lr{CQI/RI/PMI} هر \acro{UE}، بازخورد \acro{HARQ} دریافتی برای زیرقاب $n-4$،
        اندازهٔ بافرها برای هر حامل، پارامترهای کیفیت خدمات.
  \item رزرو منابع برای ارسال‌های مجدد \acro{HARQ} معوق (برای کم نگه داشتن تأخیر اولویت دارند).
  \item خدمت‌دهی به حامل‌های \acro{GBR} که زیر نرخ تضمین‌شده‌اند (مثلاً \acro{VoLTE})، با \acro{MCS} برگرفته از
        \acro{CQI} اصلاح‌شده (آفست \acro{OLLA}).
  \item رتبه‌بندی \acro{UE}های دارای داده با معیار \acro{PF} برای هر گروه \acro{RB} و تخصیص \acro{RB}ها.
  \item برای هر \acro{UE} زمان‌بندی‌شده، انتخاب \acro{MCS} و اندازهٔ بلوک انتقال و ساخت \acro{PDU}های \acro{MAC}
        (تسهیم کانال‌های منطقی بر اساس اولویت).
  \item کدگذاری \acro{DCI}ها در نامزدهای \acro{PDCCH} (اگر ناحیهٔ کنترل پر باشد، \acro{UE}ها به \acro{TTI} بعدی موکول
        می‌شوند) و ارسال \acro{PDCCH} و \acro{PDSCH}.
  \item به‌روزرسانی میانگین‌های \acro{PF} و آفست‌های \acro{OLLA} هنگام رسیدن \lr{ACK/NACK} چهار زیرقاب بعد.
\end{steps}
\end{procedurebox}

\begin{procedurebox}{دسترسی فراسو برای \lr{UE} دارای دادهٔ جدید}\label{proc:ul}
\begin{steps}
  \item دادهٔ جدید به یک بافر خالی می‌رسد؛ \acro{UE} یک \acro{BSR} راه می‌اندازد. چون مجوز ندارد، در فرصت \acro{SR}
        بعدی روی \acro{PUCCH} درخواست زمان‌بندی می‌فرستد (یا اگر پیکربندی نشده باشد دسترسی تصادفی).
  \item \acro{eNB} یک مجوز کوچک می‌فرستد (\acro{DCI} قالب ۰)؛ چهار زیرقاب بعد \acro{UE} گزارش \acro{BSR} (و \acro{PHR})
        را روی \acro{PUSCH} ارسال می‌کند.
  \item \acro{eNB} بر اساس بافر گزارش‌شده، \acro{SINR} فراسوی آن (از \lr{SRS/DMRS}) و حاشیهٔ توانش، \acro{UE} را با
        \acro{RB}های پیوسته زمان‌بندی می‌کند.
  \item هر ارسال \acro{PUSCH} پس از چهار زیرقاب روی \acro{PHICH} تأیید می‌شود؛ ارسال‌های مجدد غیرتطبیقی هشت زیرقاب
        پس از ارسال اصلی انجام می‌شوند.
\end{steps}
\end{procedurebox}

% ---------------------------------------------------------------------
\section{مثال‌های حل‌شده}

\begin{examplebox}{شمارش منابع}\label{ex:count}
\textbf{مسئله.} یک حامل \acro{LTE} پانزده‌مگاهرتزی در هر ثانیه چند عنصر منبع ارائه می‌دهد و با ناحیهٔ کنترل سه‌نمادی
و چهار درگاه \acro{CRS} چند عدد از آن‌ها عنصر منبع \acro{PDSCH} هستند؟

\textbf{حل.} $75\times12\times14\times1000 = 12.6\times10^6$ عنصر منبع در ثانیه. عناصر داده در هر جفت \acro{RB}:
$168 - 36 - 16 = 116$ (چهار درگاه \acro{CRS} از ۲۴ عنصر منبع استفاده می‌کنند که ۸ عدد آن‌ها در سه نماد کنترل
قرار دارند). مجموع $75\times116\times1000 = 8.7\times10^6$ عنصر منبع در ثانیه، یعنی ۶۹٪.
\end{examplebox}

\begin{examplebox}{گام \lr{OLLA}}\label{ex:olla}
\textbf{مسئله.} اپراتوری برای یک خدمت تحمل‌پذیر در برابر تأخیر هدف \lr{BLER} بیست درصد را می‌خواهد و از
$\delta_{\mathrm{down}} = 1$\,\lr{dB} استفاده می‌کند. $\delta_{\mathrm{up}}$ و رانش آفست را در صورتی که \lr{BLER} واقعی
۱۰٪ باشد بیابید.

\textbf{حل.} $\delta_{\mathrm{up}} = 1\times0.2/0.8 = 0.25$\,\lr{dB}. در \lr{BLER} ده درصد تغییر مورد انتظار به ازای هر
ارسال $0.9\times0.25 - 0.1\times1 = +0.125$\,\lr{dB} است: آفست بالا می‌رود (\acro{MCS} تهاجمی‌تر) تا \lr{BLER} به ۲۰٪ برسد.
\end{examplebox}

\begin{examplebox}{تصمیم \lr{PF}}\label{ex:pf}
\textbf{مسئله.} دو \acro{UE}: \lr{UE A} میانگین \lr{10\,Mbit/s} و نرخ لحظه‌ای \lr{18\,Mbit/s} دارد؛ \lr{UE B} میانگین
\lr{2\,Mbit/s} و نرخ لحظه‌ای \lr{5\,Mbit/s}. بیشینهٔ \lr{C/I} و \acro{PF} به کدام خدمت می‌دهند؟

\textbf{حل.} بیشینهٔ \lr{C/I}: \lr{UE A} (۱۸ > ۵). \acro{PF}: $18/10 = 1.8$ در برابر $5/2 = 2.5$ $\Rightarrow$ \lr{UE B} که
نسبت به میانگین خودش در قلهٔ محوشدگی است.
\end{examplebox}

% ---------------------------------------------------------------------
\section{مقایسهٔ فناوری‌ها}
\begin{table}[htbp]
\centering\small
\caption{زمان‌بندی و تطبیق پیوند در نسل‌های مختلف.}\label{tab:gens}
\rowcolors{2}{acnmist}{white}
\begin{tabularx}{\linewidth}{@{}l Y Y Y@{}}
\toprule
\rowcolor{white} & \acro{HSDPA} (\lr{3.5G}) & \acro{LTE} (\lr{4G}) & \acro{NR} (\lr{5G}) \\
\midrule
واحد زمان‌بندی    & \acro{TTI} دومیلی‌ثانیه‌ای، کدها & \acro{TTI} یک‌میلی‌ثانیه‌ای، جفت‌های \acro{RB} (\lr{180\,kHz}) & شیار (۱ تا ۰٫۱۲۵ میلی‌ثانیه) یا ریزشیار، \acro{RB}های $12\Delta f$ \\
حوزهٔ فرکانس      & ندارد (\acro{CDMA})          & به ازای گروه \acro{RB} (\acro{OFDMA}) & به ازای \acro{RB}، بخش‌های پهنای باند \\
بازخورد کانال     & \acro{CQI}                   & \lr{CQI, PMI, RI} (زیرباند)       & مبتنی بر \lr{CSI-RS}، گزارش پرتو \\
\acro{HARQ}       & \lr{NodeB}، ۶ فرایند         & ۸ فرایند، رفت‌وبرگشت ۸ میلی‌ثانیه (\acro{FDD}) & تا ۱۶ فرایند، زمان‌بندی انعطاف‌پذیر \\
فراسو             & \lr{E-DCH} (نسخهٔ ۶)          & \acro{SC-FDMA}، \acro{RB}های پیوسته & \lr{CP-OFDM} یا \lr{DFT-s-OFDM}، مجوزهای پیکربندی‌شده \\
\bottomrule
\end{tabularx}
\end{table}

\paragraph{مصالحه‌های طراحی} دانه‌بندی ریزتر تطبیق بهتر را ممکن می‌کند اما سربار کنترل دارد (\acro{DCI}های بیشتر،
گزارش‌های \acro{CQI} بیشتر)؛ بازخورد مکرر دقت را بهبود می‌دهد اما ظرفیت فراسو مصرف می‌کند؛ پنجره‌های طولانی‌تر
\acro{PF} به نفع گذردهی و پنجره‌های کوتاه‌تر به نفع انصاف کوتاه‌مدت و تأخیرند. \lr{5G} همین اصول را با عدد‌شناسی
انعطاف‌پذیر گسترش می‌دهد (درس ۱۴).

% ---------------------------------------------------------------------
\section{نکته‌های کلیدی}
\begin{takeawaybox}
\begin{itemize}
  \item قاب ۱۰ میلی‌ثانیه $\to$ زیرقاب ۱ میلی‌ثانیه $\to$ شیار ۰٫۵ میلی‌ثانیه (۷ نماد)؛ \acro{RB} $=$ ۱۲ زیرحامل
        $\times$ شیار؛ $N_{\mathrm{RB}} = 6\dots100$ برای ۱٫۴ تا ۲۰ مگاهرتز.
  \item \acro{3GPP} سیگنالینگ (\lr{DCI, CQI, BSR, PHR}، زمان‌بندی \acro{HARQ}) را استاندارد می‌کند، نه الگوریتم زمان‌بندی را.
  \item \acro{CQI} $\to$ بازده $Q_m\times$ نرخ کد؛ \acro{OLLA} مقدار \lr{BLER} را به $\delta_{\mathrm{up}}/(\delta_{\mathrm{up}}+\delta_{\mathrm{down}})$ می‌راند.
  \item هشت فرایند \acro{HARQ} کانال را مشغول نگه می‌دارند؛ \lr{BLER} ده درصد ارسال نخست حدود ۹٪ هزینه دارد.
  \item بیشینهٔ \lr{C/I} مقدار $\sum \bar R_k$ را، \acro{PF} مقدار $\sum\log\bar R_k$ را بیشینه می‌کند و \acro{RR} زمان را برابر می‌کند.
  \item تنوع چندکاربره: در کانال‌های محوشونده \acro{PF} می‌تواند برای \emph{هر} کاربر از \acro{RR} بهتر باشد.
  \item زمان‌بندهای عملی: ابتدا ارسال‌های مجدد \acro{HARQ}، سپس \acro{GBR}/کیفیت خدمات و سپس \acro{PF} برای بهترین-تلاش.
\end{itemize}
\end{takeawaybox}

% ---------------------------------------------------------------------
\section{تمرین‌ها}

\begin{exercisebox}{1}{استاندارد در برابر سازنده}\label{exr:std}
کدام یک از موارد زیر را \acro{3GPP} استاندارد کرده و کدام پیاده‌سازی سازنده است؟ جدول \acro{CQI}، معیار منصفانهٔ
تناسبی، قالب \acro{DCI}، اندازهٔ گام‌های \acro{OLLA}، زمان‌بندی \acro{HARQ} در \acro{FDD}، اولویت حامل‌های \acro{GBR}.
\end{exercisebox}

\begin{exercisebox}{1}{پیوستگی فراسو}\label{exr:contig}
توضیح دهید چرا \acro{SC-FDMA} برای هر \acro{UE} به \acro{RB}های پیوسته نیاز دارد و این برای زمان‌بندی فراسوی گزینشی
در فرکانس چه معنایی دارد.
\end{exercisebox}

\begin{exercisebox}{2}{نرخ از روی \lr{CQI}}\label{exr:cqi}
یک \acro{UE} روی حامل بیست‌مگاهرتزی (یک نماد کنترل، دو درگاه \acro{CRS}) مقدار \acro{CQI} ۱۲ با رتبهٔ ۲ گزارش
می‌دهد. ۴۰ \acro{RB} دریافت می‌کند. نرخ آن را با \cref{eq:cqirate} محاسبه کنید.
\end{exercisebox}

\begin{exercisebox}{2}{شاخص \lr{Jain}}\label{exr:jain}
شاخص‌های \lr{Jain} \cref{tab:sim} را برای \acro{RR} و بیشینهٔ \lr{C/I} از روی مقادیر هر کاربر بررسی کنید.
\end{exercisebox}

\begin{exercisebox}{3}{تحلیل \lr{PF} دوکاربره}\label{exr:pf2}
دو کاربر کانال‌های هم‌توزیع دارند اما نرخ کاربر ۱ همیشه دو برابر کاربر ۲ است ($r_1 = 2r_2$ در هر \acro{TTI}). نشان
دهید \acro{PF} به هر کاربر نیمی از \acro{TTI}ها را می‌دهد و گذردهی‌های حاصل را نسبت به \acro{RR} محاسبه کنید. اگر
کانال‌ها مستقل باشند چه تغییری رخ می‌دهد؟
\end{exercisebox}

\begin{exercisebox}{3}{مطالعهٔ شبیه‌سازی}\label{exr:sim}
\lr{\texttt{scheduler\_sim.py}} را برای ۸ کاربر با میانگین \acro{SNR} یکنواخت میان ۰ تا ۲۱ دسی‌بل تغییر دهید و
گذردهی سلول و شاخص \lr{Jain} را بر حسب $t_\mathrm{c}\in\{10,100,1000\}$ برای \acro{PF} رسم کنید. روند را توضیح دهید.
\end{exercisebox}

\subsection*{پاسخ‌ها}
\addcontentsline{toc}{subsection}{پاسخ‌ها}

\begin{solutionbox}{\ref{exr:std}}
استانداردشده: جدول \acro{CQI}، قالب‌های \acro{DCI}، زمان‌بندی \acro{HARQ} در \acro{FDD} (و معنای \lr{BSR/PHR}). پیاده‌سازی
سازنده: معیار \acro{PF} (یا هر الگوریتمی)، اندازهٔ گام‌های \acro{OLLA}، سیاست اولویت‌بندی میان حامل‌ها (که پارامترهای
کیفیت خدمات آن را محدود می‌کنند اما با جزئیات تعیین نمی‌کنند).
\end{solutionbox}

\begin{solutionbox}{\ref{exr:contig}}
\acro{SC-FDMA} همان \acro{OFDM} با گسترش \acro{DFT} است: نگاشت خروجی \acro{DFT} به زیرحامل‌های پیوسته سیگنالی شبیه
تک‌حامل با نسبت توان قله به میانگین پایین تولید می‌کند و تقویت‌کننده‌های توان کارآمد را در گوشی ممکن می‌سازد.
نگاشت غیرپیوسته این نسبت را بالا می‌برد. بنابراین زمان‌بند می‌تواند انتخاب کند بلوک یک \acro{UE} \emph{کجا} قرار
گیرد (بهرهٔ گزینشی در فرکانس) اما نمی‌تواند \acro{RB}های خوب پراکنده را برگزیند؛ نسخهٔ ۱۰ این را با تخصیص
خوشه‌ای دو خوشه‌ای تعدیل کرد.
\end{solutionbox}

\begin{solutionbox}{\ref{exr:cqi}}
عناصر منبع داده در هر جفت \acro{RB}: $168 - 12 - (16-4) = 144$. $R = 2\times3.9023\times40\times144/10^{-3} =
\qty{45.0}{\mega\bit\per\second}$.
\end{solutionbox}

\begin{solutionbox}{\ref{exr:jain}}
\acro{RR}: مجموع $=19.88$، مجموع مربعات $=94.87+33.99+8.29+2.04 = 139.2$؛ $J = 395.2/(4\times139.2) = 0.710$.
بیشینهٔ \lr{C/I}: مجموع $= 40.13$، مربعات $\approx 1276.6+17.5+0.04 = 1294.2$؛ $J = 1610.4/(4\times1294.2) = 0.311$.
\end{solutionbox}

\begin{solutionbox}{\ref{exr:pf2}}
با $r_1 = 2r_2$ در همهٔ زمان‌ها، مقایسهٔ $r_1/\bar R_1$ با $r_2/\bar R_2$ یعنی مقایسهٔ $2r_2/\bar R_1$ با $r_2/\bar R_2$؛
تعادل مستلزم انتخاب هر دو است که وقتی رخ می‌دهد که $\bar R_1 = 2\bar R_2$، یعنی هر کاربر نیمی از \acro{TTI}ها را
بگیرد -- دقیقاً \acro{RR}، زیرا کانال‌ها کاملاً همبسته‌اند و تنوع چندکاربره‌ای برای بهره‌برداری وجود ندارد. با
کانال‌های مستقل، \acro{PF} هر کاربر را در قله‌های نسبی‌اش انتخاب می‌کند و هر دو گذردهی از \acro{RR} بیشتر می‌شوند،
مانند \cref{tab:sim}.
\end{solutionbox}

\begin{solutionbox}{\ref{exr:sim}}
با میانگین \acro{SNR} برابر $0,3,\dots,21$ دسی‌بل و ۴۰۰۰۰ \acro{TTI} اسکریپت این نتایج را می‌دهد:
$t_\mathrm{c}=10$: سلول \lr{29.8\,Mbit/s} و $J=0.831$؛ $t_\mathrm{c}=100$: \lr{32.8\,Mbit/s} و $J = 0.838$؛
$t_\mathrm{c}=1000$: \lr{33.2\,Mbit/s} و $J=0.839$ (\acro{RR}: \lr{21.8\,Mbit/s} و $J=0.775$). پنجرهٔ کوتاه به خدمت اخیر
واکنش نشان می‌دهد و گاهی به کاربری خدمت می‌دهد که در قلهٔ نسبی نیست؛ پنجره‌های طولانی‌تر به \acro{PF} اجازه
می‌دهند منتظر قله‌های واقعی بماند؛ پس گذردهی بالا می‌رود و وقتی پنجره بسیار طولانی‌تر از محوشدگی شود (که اینجا
از \acro{TTI} به \acro{TTI} مستقل است) اشباع می‌شود. بهای پنجره‌های طولانی در عمل، بی‌انصافی کوتاه‌مدت و لرزش تأخیر
بیشتر در کانال‌های همبستهٔ زمانی است.
\end{solutionbox}

% ---------------------------------------------------------------------
\section{مطالعهٔ بیشتر}
کانال‌های فیزیکی و شبکهٔ منابع در \cite{ts36211}، رویه‌های \acro{CQI}، \acro{DCI} و \acro{HARQ} در \cite{ts36213} و
رویه‌های \acro{MAC} (\lr{BSR, PHR, SR}) در \cite{ts36321} مشخص شده‌اند. زمان‌بند منصفانهٔ تناسبی از \lr{CDMA-HDR}
سرچشمه گرفت \cite{jalali2000hdr}؛ تفسیر تنوع چندکاربرهٔ آن در \lr{Tse} و \lr{Viswanath} \cite{tse2005fundamentals}
بسط داده شده است. شاخص \lr{Jain} از \cite{jain1984fairness} است. \lr{Sesia} و همکاران \cite{sesia2011lte} زمان‌بندی
\acro{LTE} را به‌تفصیل پوشش می‌دهند.

\begin{latin}
\printbibliography[heading=subbibliography,title={References}]
\end{latin}
\end{document}
